\documentclass[11pt,a4paper]{article}
\usepackage[utf8]{inputenc}
\usepackage[T1]{fontenc}
\usepackage{geometry}
\geometry{margin=1in}
\usepackage{hyperref}
\usepackage{enumitem}
\usepackage{booktabs}
\usepackage{longtable}
\usepackage{xcolor}
\usepackage{titlesec}
\usepackage{amsmath,amssymb}

\titleformat{\section}{\Large\bfseries\color{blue!70!black}}{}{0em}{}
\titleformat{\subsection}{\large\bfseries\color{blue!50!black}}{}{0em}{}
\titleformat{\subsubsection}{\normalsize\bfseries\color{blue!30!black}}{}{0em}{}

\title{\textbf{Replication Plan}\\[0.5em]
\large Latent Stochastic Differential Equations for Modeling Quasar Variability and Inferring Black Hole Properties\\[0.3em]
\normalsize OSTI ID: 2396968}
\author{Auto-generated Replication Blueprint}
\date{\today}

\begin{document}
\maketitle
\tableofcontents
\newpage

%---------------------------------------------------------------------
\section{Paper Overview}
%---------------------------------------------------------------------
This paper develops latent stochastic differential equation (SDE) models for quasar optical variability and uses them to infer physical properties of supermassive black holes (mass, accretion rate). Quasar light curves are modeled as observations of a latent SDE process, going beyond the standard damped random walk (DRW/OU process) to capture more complex variability. The model is fit using likelihood-based inference (possibly with the \texttt{celerite} fast GP framework for scalable computation). The authors validate on simulated light curves and apply to real survey data (e.g., SDSS Stripe 82, ZTF).

%---------------------------------------------------------------------
\section{Detailed Workflow}
%---------------------------------------------------------------------

\subsection{Phase 1: Stochastic Process Models}
\begin{enumerate}[label=\textbf{1.\arabic*}]
  \item \textbf{Implement the DRW (OU process) baseline.}
    \begin{itemize}
      \item $dX(t) = -\frac{1}{\tau} X(t)\,dt + \sigma\,dW(t)$
      \item Analytic covariance function: $C(\Delta t) = \hat{\sigma}^2 e^{-|\Delta t|/\tau}$.
    \end{itemize}
  \item \textbf{Implement higher-order latent SDEs.}
    \begin{itemize}
      \item CARMA($p,q$) models or custom multi-dimensional SDEs.
      \item These have richer power spectral densities (PSDs).
    \end{itemize}
  \item \textbf{Implement the celerite representation} for fast ($O(N)$) likelihood evaluation.
\end{enumerate}

\subsection{Phase 2: Simulated Light Curve Generation}
\begin{enumerate}[label=\textbf{2.\arabic*}]
  \item \textbf{Generate mock light curves} from each SDE model.
    \begin{itemize}
      \item Sample at irregular time cadences mimicking survey strategies.
      \item Add photometric noise with realistic magnitude-dependent errors.
    \end{itemize}
  \item \textbf{Inject known physical parameters} ($M_{\text{BH}}$, $\dot{M}$, $\tau$, $\sigma$).
\end{enumerate}

\subsection{Phase 3: Parameter Inference}
\begin{enumerate}[label=\textbf{3.\arabic*}]
  \item \textbf{Likelihood computation} using celerite or direct GP.
  \item \textbf{Parameter optimization} (maximum likelihood) and \textbf{MCMC sampling} (posterior distributions).
    \begin{itemize}
      \item Use emcee or PyMC for MCMC.
      \item Sample SDE parameters and physical parameter mappings.
    \end{itemize}
  \item \textbf{Model comparison} (DRW vs.\ higher-order SDE) via BIC, AIC, or Bayes factors.
\end{enumerate}

\subsection{Phase 4: Application to Real Data}
\begin{enumerate}[label=\textbf{4.\arabic*}]
  \item \textbf{Download quasar light curves} from SDSS Stripe 82 or ZTF.
  \item \textbf{Fit models} to real light curves.
  \item \textbf{Infer black hole masses} and compare to spectroscopic estimates.
  \item \textbf{Population-level analysis} of inferred parameters.
\end{enumerate}

\subsection{Phase 5: Validation}
\begin{enumerate}[label=\textbf{5.\arabic*}]
  \item Parameter recovery tests on simulated data.
  \item Reproduce paper's figures (posterior distributions, PSD comparisons, mass correlations).
\end{enumerate}

%---------------------------------------------------------------------
\section{Datasets}
%---------------------------------------------------------------------

\begin{longtable}{p{4.5cm} p{3.5cm} p{6cm}}
\toprule
\textbf{Dataset / Source} & \textbf{Type} & \textbf{Location / Notes} \\
\midrule
\endfirsthead
\toprule
\textbf{Dataset / Source} & \textbf{Type} & \textbf{Location / Notes} \\
\midrule
\endhead
Simulated quasar light curves & Synthetic & Generated from SDE models with known parameters \\
\midrule
SDSS Stripe 82 light curves & Public survey & \url{https://www.sdss.org/dr17/} \newline Quasar catalog: \url{https://faculty.washington.edu/ivezic/macleod/qso_dr7/} \\
\midrule
ZTF light curves & Public survey & \url{https://www.ztf.caltech.edu/} \\
\midrule
Quasar catalog (spectroscopic) & Public & SDSS DR16Q: \url{https://www.sdss.org/dr16/algorithms/qso_catalog/} \\
\midrule
Black hole mass estimates & Published & Shen et al.\ (2011) SDSS quasar catalog \\
\bottomrule
\end{longtable}

%---------------------------------------------------------------------
\section{Models, Tools, and Software}
%---------------------------------------------------------------------

\begin{longtable}{p{4.5cm} p{2.5cm} p{6.5cm}}
\toprule
\textbf{Tool / Library} & \textbf{Version} & \textbf{Purpose / URL} \\
\midrule
\endfirsthead
\toprule
\textbf{Tool / Library} & \textbf{Version} & \textbf{Purpose / URL} \\
\midrule
\endhead
celerite / celerite2 & latest & Fast GP/SDE likelihood computation \newline \url{https://celerite2.readthedocs.io/} \\
\midrule
emcee & $\geq 3.1$ & MCMC sampler \newline \url{https://emcee.readthedocs.io/} \\
\midrule
PyMC & $\geq 5.0$ & Bayesian modeling framework (alternative) \newline \url{https://www.pymc.io/} \\
\midrule
JAX & latest & Autodiff and JIT compilation for likelihood \newline \url{https://github.com/google/jax} \\
\midrule
NumPy / SciPy & $\geq 1.21$ & Numerical operations \\
\midrule
Matplotlib & $\geq 3.4$ & Light curve and posterior plots \\
\midrule
Astropy & $\geq 5.0$ & Astronomical data I/O and coordinates \\
\midrule
corner.py & latest & Corner plots for MCMC posteriors \newline \url{https://corner.readthedocs.io/} \\
\bottomrule
\end{longtable}

\subsection{Hardware Requirements}
\begin{itemize}
  \item MCMC on individual light curves: single CPU core, minutes per object.
  \item Population analysis ($\sim$10,000 quasars): multi-core workstation or cluster.
\end{itemize}

\end{document}
